% !TeX encoding = UTF-8
\documentclass[variant=guia,cover=institutional,mode=screen]{ua-engineering-v6-article}
% !TeX encoding = UTF-8
\UASetUniversity{Universidad Austral}
\UASetFaculty{Facultad de Ingeniería}
\UASetDepartment{Departamento de Computación}
\UASetCampus{Pilar}
\UASetCareer{Ingeniería Informática}
\UASetCommission{Comisión A}
\UASetProfessor{Equipo docente}
\UASetSemester{Segundo cuatrimestre 2026}
\UASetCourse{Sistemas Distribuidos}
\UASetDate{2026}


\author{\UAProfessor}
\date{\UADate}
\begin{document}

\section{Indicaciones generales}
Resuelva cada ejercicio declarando supuestos, modelo de fallas, garantía y trade-off. Cuando corresponda, construya una ejecución verificable y vincule la respuesta con evidencia observable.

\section{Nivel 1: fundamentos y reconocimiento}
\begin{uaexercise}
Defina atomicidad, consistencia, aislamiento y durabilidad, y explique por qué cada propiedad responde a un problema distinto.
\end{uaexercise}

\begin{uaexercise}
Diferencie atomicidad de aislamiento mediante una ejecución donde una se preserve y la otra no.
\end{uaexercise}

\begin{uaexercise}
Dado un schedule con dos transacciones, identifique conflictos RW, WR y WW.
\end{uaexercise}

\begin{uaexercise}
Construya el grafo de precedencia de un schedule y determine si es conflict-serializable.
\end{uaexercise}

\begin{uaexercise}
Demuestre informalmente por qué un ciclo en el grafo impide un orden serial equivalente.
\end{uaexercise}


\section{Nivel 2: ejecuciones y fallas}
\begin{uaexercise}
Construya un dirty read y explique el estado observable después de un abort.
\end{uaexercise}

\begin{uaexercise}
Construya una non-repeatable read e indique qué nivel de aislamiento la permite.
\end{uaexercise}

\begin{uaexercise}
Construya un phantom read con una consulta por predicado.
\end{uaexercise}

\begin{uaexercise}
Desarrolle el caso de write skew de dos médicos y formule el invariante violado.
\end{uaexercise}

\begin{uaexercise}
Explique por qué Snapshot Isolation evita conflictos WW pero puede permitir write skew.
\end{uaexercise}


\section{Nivel 3: diseño y comparación}
\begin{uaexercise}
Compare Read Committed, Repeatable Read, Snapshot Isolation y Serializable usando anomalías permitidas.
\end{uaexercise}

\begin{uaexercise}
Diseñe una prueba que diferencie Snapshot Isolation de serializabilidad real.
\end{uaexercise}

\begin{uaexercise}
Describa el protocolo 2PC con coordinador y tres participantes, incluyendo registros durables.
\end{uaexercise}

\begin{uaexercise}
Analice una caída del coordinador antes de prepare, durante prepare y después de recibir todos los votos YES.
\end{uaexercise}

\begin{uaexercise}
Explique por qué un participante prepared no puede decidir unilateralmente sin riesgo para atomicidad.
\end{uaexercise}


\section{Nivel 4: integración y defensa}
\begin{uaexercise}
Compare 2PC con consenso: problema que resuelve cada uno, seguridad y progreso.
\end{uaexercise}

\begin{uaexercise}
Explique qué supuesto adicional necesita 3PC y por qué no elimina las particiones reales.
\end{uaexercise}

\begin{uaexercise}
Diseñe una saga para reserva, pago y emisión; incluya compensaciones e idempotencia.
\end{uaexercise}

\begin{uaexercise}
Aplique P-F-G-S-T a una transacción distribuida de su Trabajo Final.
\end{uaexercise}

\begin{uaexercise}
Defienda ante comité cuándo elegiría transacción fuerte, 2PC o saga para una operación crítica.
\end{uaexercise}


\end{document}
